\documentclass[12pt,a4paper]{ctexart}

% ===== 页面设置 =====
\usepackage[top=2.5cm,bottom=2.5cm,left=2cm,right=2cm]{geometry}
\usepackage{booktabs}
\usepackage{longtable}
\usepackage{array}
\usepackage{xcolor}
\usepackage{colortbl}
\usepackage{hyperref}

\hypersetup{colorlinks=true,linkcolor=blue}

\title{\textbf{办公物品管理系统\\数据字典与数据结构}}
\author{}
\date{\today}

\begin{document}
\maketitle
\tableofcontents
\newpage

\section{表汇总}

\begin{table}[h]
\centering
\begin{tabular}{c|c|c|c|c}
\toprule
\textbf{序号} & \textbf{表名} & \textbf{中文名} & \textbf{引擎} & \textbf{记录数（示例）} \\
\midrule
1 & Categories & 物品类别字典 & InnoDB & 6 \\
2 & Origins & 产地字典 & InnoDB & 13 \\
3 & Items & 物品信息表 & InnoDB & 10 \\
4 & StockIn & 入库记录表 & InnoDB & $\sim$1,000,010 \\
5 & StockOut & 领用出库记录表 & InnoDB & 3 \\
\bottomrule
\end{tabular}
\caption{数据库表汇总}
\end{table}

\section{表结构详细定义}

\subsection{Categories（物品类别字典）}

\begin{longtable}{c|c|c|c|c}
\toprule
\textbf{字段名} & \textbf{数据类型} & \textbf{允许空} & \textbf{主键/外键} & \textbf{说明} \\
\midrule
\endfirsthead
\toprule
\textbf{字段名} & \textbf{数据类型} & \textbf{允许空} & \textbf{主键/外键} & \textbf{说明} \\
\midrule
\endhead
Id   & INT & 否 & PK\_Categories & 类别 ID，自增 \\
Name & VARCHAR(50) & 否 & UQ\_Categories\_Name & 类别名称，唯一 \\
\bottomrule
\caption{Categories 表结构}
\end{longtable}

\vspace{4pt}
\textbf{预置数据：}

\begin{tabular}{c|c}
\toprule
\textbf{Id} & \textbf{Name} \\
\midrule
1 & 纸张 \\
2 & 文具 \\
3 & 刀具 \\
4 & 单据 \\
5 & 礼品 \\
6 & 其它 \\
\bottomrule
\end{tabular}

\subsection{Origins（产地字典）}

\begin{longtable}{c|c|c|c|c}
\toprule
\textbf{字段名} & \textbf{数据类型} & \textbf{允许空} & \textbf{主键/外键} & \textbf{说明} \\
\midrule
\endfirsthead
\toprule
\textbf{字段名} & \textbf{数据类型} & \textbf{允许空} & \textbf{主键/外键} & \textbf{说明} \\
\midrule
\endhead
Id   & INT & 否 & PK\_Origins & 产地 ID，自增 \\
Name & VARCHAR(100) & 否 & UQ\_Origins\_Name & 产地名称，唯一 \\
\bottomrule
\caption{Origins 表结构}
\end{longtable}

\vspace{4pt}
\textbf{预置数据：} 北京、上海、广州、深圳、杭州、成都、武汉、南京、苏州、德国、日本、美国、中国台湾（共 13 个）

\subsection{Items（物品信息表）}

\begin{longtable}{c|c|c|c|c}
\toprule
\textbf{字段名} & \textbf{数据类型} & \textbf{允许空} & \textbf{主键/外键} & \textbf{说明} \\
\midrule
\endfirsthead
\toprule
\textbf{字段名} & \textbf{数据类型} & \textbf{允许空} & \textbf{主键/外键} & \textbf{说明} \\
\midrule
\endhead
ItemCode      & VARCHAR(20)  & 否 & PK\_Items            & 物品编码，唯一标识 \\
ItemName      & VARCHAR(100) & 否 & —                    & 物品名称 \\
Category      & INT          & 否 & FK → Categories(Id)  & 物品类别 \\
Origin        & VARCHAR(100) & 否（默认''）& —                & 产地 \\
Specification & VARCHAR(100) & 否（默认''）& —               & 规格描述 \\
Model         & VARCHAR(100) & 否（默认''）& —              & 型号 \\
ImagePath     & VARCHAR(500) & 否（默认''）& —             & 图片路径（兼容历史） \\
ImageData     & LONGBLOB     & 是 & —                    & 图片二进制数据 \\
Quantity      & INT          & 否（默认 0）& CK\_Items\_Quantity & 库存数量，CHECK($\geq$0) \\
\bottomrule
\caption{Items 表结构}
\end{longtable}

\vspace{4pt}
\textbf{索引：}

\begin{tabular}{c|c|c}
\toprule
\textbf{索引名} & \textbf{字段} & \textbf{说明} \\
\midrule
IX\_Items\_Category & Category  & 加速按类别筛选 \\
IX\_Items\_ItemName & ItemName  & 加速按名称搜索 \\
\bottomrule
\end{tabular}

\subsection{StockIn（入库记录表）}

\begin{longtable}{c|c|c|c|c}
\toprule
\textbf{字段名} & \textbf{数据类型} & \textbf{允许空} & \textbf{主键/外键} & \textbf{说明} \\
\midrule
\endfirsthead
\toprule
\textbf{字段名} & \textbf{数据类型} & \textbf{允许空} & \textbf{主键/外键} & \textbf{说明} \\
\midrule
\endhead
StockInNo    & VARCHAR(20)    & 否 & PK\_StockIn              & 流水号 SI+日期+序号 \\
ItemCode     & VARCHAR(20)    & 否 & FK → Items(ItemCode)     & 物品编码 \\
PurchaseDate & DATETIME       & 否 & —                        & 购买日期 \\
Quantity     & INT            & 否 & CK\_StockIn\_Quantity     & 入库数量，CHECK($>$0) \\
UnitPrice    & DECIMAL(18,2)  & 否 & CK\_StockIn\_UnitPrice    & 单价，CHECK($\geq$0) \\
TotalPrice   & DECIMAL(18,2)  & 否 & CK\_StockIn\_TotalPrice   & 总价=数量×单价 \\
\bottomrule
\caption{StockIn 表结构}
\end{longtable}

\vspace{4pt}
\textbf{索引：}

\begin{tabular}{c|c|c}
\toprule
\textbf{索引名} & \textbf{字段} & \textbf{说明} \\
\midrule
IX\_StockIn\_ItemCode     & ItemCode      & 加速按物品查入库 \\
IX\_StockIn\_PurchaseDate & PurchaseDate  & 加速按日期范围查询 \\
\bottomrule
\end{tabular}

\subsection{StockOut（领用出库记录表）}

\begin{longtable}{c|c|c|c|c}
\toprule
\textbf{字段名} & \textbf{数据类型} & \textbf{允许空} & \textbf{主键/外键} & \textbf{说明} \\
\midrule
\endfirsthead
\toprule
\textbf{字段名} & \textbf{数据类型} & \textbf{允许空} & \textbf{主键/外键} & \textbf{说明} \\
\midrule
\endhead
StockOutNo  & VARCHAR(20)   & 否 & PK\_StockOut               & 流水号 SO+日期+序号 \\
ItemCode    & VARCHAR(20)   & 否 & FK → Items(ItemCode)       & 物品编码 \\
Quantity    & INT           & 否 & CK\_StockOut\_Quantity     & 领用数量，CHECK($>$0) \\
ApplicantId & VARCHAR(50)   & 否 & —                          & 领用人 ID（JSON 用户） \\
ApplyDate   & DATETIME      & 否 & —                          & 申请日期 \\
Status      & INT           & 否（默认 0）& CK\_StockOut\_Status   & 0=申请 1=确认 2=驳回 \\
ApproveDate & DATETIME      & 是 & —                          & 审批日期 \\
Remark      & VARCHAR(500)  & 是 & —                          & 审批备注 \\
\bottomrule
\caption{StockOut 表结构}
\end{longtable}

\vspace{4pt}
\textbf{索引：}

\begin{tabular}{c|c|c}
\toprule
\textbf{索引名} & \textbf{字段} & \textbf{说明} \\
\midrule
IX\_StockOut\_ItemCode    & ItemCode    & 加速按物品查询 \\
IX\_StockOut\_ApplicantId & ApplicantId & 加速按用户查询 \\
IX\_StockOut\_Status      & Status      & 加速按状态筛选 \\
\bottomrule
\end{tabular}

\section{状态码定义}

\subsection{StockOut.Status（领用状态）}

\begin{table}[h]
\centering
\begin{tabular}{c|c|c|c}
\toprule
\textbf{值} & \textbf{含义} & \textbf{触发者} & \textbf{库存影响} \\
\midrule
0 & 申请 & 普通用户 & 不变 \\
1 & 确认 & 管理员   & \cellcolor{red!20}\textbf{扣减} \\
2 & 驳回 & 管理员   & 不变 \\
\bottomrule
\end{tabular}
\caption{领用状态码定义}
\end{table}

\subsection{Items.Category（物品类别）}

\begin{table}[h]
\centering
\begin{tabular}{c|c}
\toprule
\textbf{值} & \textbf{含义} \\
\midrule
1 & 纸张 \\
2 & 文具 \\
3 & 刀具 \\
4 & 单据 \\
5 & 礼品 \\
6 & 其它 \\
\bottomrule
\end{tabular}
\caption{物品类别码}
\end{table}

\section{流水号编码规则}

\begin{longtable}{c|c|c|c}
\toprule
\textbf{表} & \textbf{格式} & \textbf{示例} & \textbf{说明} \\
\midrule
\endfirsthead
\toprule
\textbf{表} & \textbf{格式} & \textbf{示例} & \textbf{说明} \\
\midrule
\endhead
StockIn  & SI + yyyyMMdd + 3位序号 & SI20250115001 & 2025年1月15日第1笔入库 \\
StockOut & SO + yyyyMMdd + 3位序号 & SO20250115002 & 2025年1月15日第2笔领用 \\
\bottomrule
\caption{流水号编码规则}
\end{longtable}

\section{视图}

\begin{longtable}{c|c|c}
\toprule
\textbf{视图名} & \textbf{数据源} & \textbf{用途} \\
\midrule
\endfirsthead
\toprule
\textbf{视图名} & \textbf{数据源} & \textbf{用途} \\
\midrule
\endhead
v\_StockSummary     & Items JOIN Categories                    & 库存总览（含库存等级） \\
v\_StockInMonthly   & StockIn GROUP BY                          & 入库按月统计 \\
v\_StockOutSummary  & StockOut JOIN Items JOIN Categories       & 领用按物品统计 \\
v\_ItemFullInfo     & Items + 子查询（入库总额、领用次数）       & 物品完整信息 \\
\bottomrule
\caption{视图清单}
\end{longtable}

\section{触发器}

\begin{longtable}{c|c|c|c}
\toprule
\textbf{触发器名} & \textbf{表} & \textbf{时机} & \textbf{操作} \\
\midrule
\endfirsthead
\toprule
\textbf{触发器名} & \textbf{表} & \textbf{时机} & \textbf{操作} \\
\midrule
\endhead
trg\_StockIn\_AfterInsert  & StockIn  & AFTER INSERT  & Items.Quantity += NEW.Quantity \\
trg\_StockIn\_AfterDelete  & StockIn  & AFTER DELETE  & Items.Quantity -= OLD.Quantity（至少为 0） \\
trg\_StockOut\_AfterUpdate & StockOut & AFTER UPDATE  & 状态 0→1 时 Items.Quantity -= NEW.Quantity \\
trg\_StockOut\_AfterDelete & StockOut & AFTER DELETE  & 已确认记录删除时恢复库存 \\
\bottomrule
\caption{触发器清单}
\end{longtable}

\section{存储过程}

\begin{longtable}{c|c|c}
\toprule
\textbf{存储过程} & \textbf{参数} & \textbf{功能} \\
\midrule
\endfirsthead
\hline
\textbf{存储过程} & \textbf{参数} & \textbf{功能} \\
\hline
\endhead
sp\_GenerateStockInNo    & OUT outNo & 生成入库流水号 \\
sp\_GenerateStockOutNo   & OUT outNo & 生成领用流水号 \\
sp\_StockIn              & IN 4 个参数 & 事务：写入库 + 更新库存 \\
sp\_ApproveStockOut      & IN 3 个参数 & 事务：确认 + 扣库存 \\
sp\_RejectStockOut       & IN 2 个参数 & 驳回领用 \\
sp\_StockStatistics      & —          & 按类别统计库存 \\
sp\_StockInDetail        & IN 3 个参数 & 入库明细分页 \\
sp\_StockOutDetail       & IN 3 个参数 & 领用明细分页 \\
sp\_FixNegativeStock     & —          & 游标批量修正负库存 \\
sp\_TopCategoryByStockIn & IN TopN    & 入库金额排行 \\
\bottomrule
\caption{存储过程清单}
\end{longtable}

\section{用户数据结构（JSON）}

存储文件：\texttt{Users.json}

\begin{verbatim}
[
  {
    "UserId": "user01",
    "Password": "1234",
    "Name": "张三",
    "Gender": "男",
    "BirthDate": "1990-05-15T00:00:00",
    "Phone": "13800138001"
  }
]
\end{verbatim}

\begin{longtable}{c|c|c}
\toprule
\textbf{字段} & \textbf{类型} & \textbf{说明} \\
\midrule
\endfirsthead
\toprule
\textbf{字段} & \textbf{类型} & \textbf{说明} \\
\midrule
\endhead
UserId    & string   & 用户 ID（登录账号） \\
Password  & string   & 登录密码 \\
Name      & string   & 真实姓名 \\
Gender    & string   & 性别："男" / "女" \\
BirthDate & DateTime & 出生日期 \\
Phone     & string   & 联系电话 \\
\bottomrule
\caption{JSON 用户数据结构}
\end{longtable}

\end{document}
